%------------------------------------------------------------------------------%
\begin{question}
  Encontre os valores máximo e mínimo absolutos de
  \[
    f(x, y) = 2 + 2x + 2y - x^2 - y^2
  \]
  na região triangular limitada pelas retas
  \(x=0\), \(y=0\) e \(y = 9-x\)
\end{question}

%------------------------------------------------------------------------------%
\begin{resolution}{Pontos críticos interiores}
  \[
    f_x(x, y)
    \pause
    = \D{}{x}\left(2 + 2x + 2y - x^2 - y^2\right)
    \pause
    = 2-2x
  \]
  \[
    \pause
    f_y(x, y)
    \pause
    = \D{}{y}\left(2 + 2x + 2y - x^2 - y^2\right)
    \pause
    = 2-2y
  \]
  \vspace{-0.5\baselineskip}
  \begin{columns}
  \column{0.3\linewidth}
    \pause
    \begin{align*}
      f_x(x, y) & = 0 \\
      2-2x      & = 0 \\
      x         & = 1
    \end{align*}
  \column{0.3\linewidth}
    \pause
    \begin{align*}
      f_y(x, y) & = 0 \\
      2-2y      & = 0 \\
      y         & = 1
    \end{align*}
  \end{columns}

  \pause
  \vspace{0.5\baselineskip}
  O ponto $(1, 1)$ está no interior da região $R$
\end{resolution}

%------------------------------------------------------------------------------%
\begin{resolution}{Valor da função}
 \[
  f(1, 1)
    \pause
    = \left(2 + 2x + 2y - x^2 - y^2\right)\evalat{(1, 1)}{}
    \pause
    = 2 + 2 + 2 - 1 - 1
    \pause
    = 4
  \]
\end{resolution}

%------------------------------------------------------------------------------%
\begin{resolution}{Pontos na fronteira}
  Vamos tratar cada segmento da fronteira separadamente
  \begin{enumerate}[<+->]
    \item \(y=0\)   \tabto{24mm} \(0\leq x\leq 9\)
    \item \(x=0\)   \tabto{24mm} \(0\leq y\leq 9\)
    \item \(y=9-x\) \tabto{24mm} \(0\leq x\leq 9\)
  \end{enumerate}
\end{resolution}

%------------------------------------------------------------------------------%
\begin{resolution}{Primeiro segmento}
  Segmento: \(y=0\), \(0\leq x\leq 9\)

  \pause
  Função de uma variável, $x$, no intervalo fechado \(0\leq x\leq 9\)
  \[
    \pause
    g(x)
    \pause
    = f(x, 0)
    \pause
    = \left(2 + 2x + 2y - x^2 - y^2\right)\evalat{y=0}{}
    \pause
    = 2 + 2x - x^2
  \]
  \pause
  Pontos críticos de $g$ no interior do intervalo \([0, 9]\)
  \vspace{-0.7\baselineskip}
  \begin{columns}[t]
  \column{0.4\linewidth}
    \pause
    \begin{align*}
      \onslide<+->{g'(x)}
      \onslide<+->{& = \frac{d}{dx}\left(2 + 2x - x^2\right) \\}
      \onslide<+->{& = 2 - 2x}
    \end{align*}
  \column{0.4\linewidth}
   \begin{align*}
      \onslide<+->{g'(x)  & = 0 \\}
      \onslide<+->{2 - 2x & = 0 \\}
      \onslide<+->{x      & = 1}
    \end{align*}
  \end{columns}
\end{resolution}

%------------------------------------------------------------------------------%
\begin{resolution}{Valores da função}
  Ponto crítico e extremos do intervalo \(x=0\), \(x=1\) e \(x=9\)

  \pause
  Pontos no plano
  \((0, 0)\), \((1, 0)\), \((9, 0)\)

  \pause
  Valores da função
  \[
    \pause f(0, 0)
    \pause = g(0)
    \pause = \left(2 + 2x - x^2\right) \evalat{0}{}
    \pause = 2
  \]
  \[
    \pause f(1, 0)
    \pause = g(1)
    \pause = \left(2 + 2x - x^2\right) \evalat{1}{}
    \pause = 3
  \]
  \[
    \pause f(9, 0)
    \pause = g(9)
    \pause = \left(2 + 2x - x^2\right) \evalat{9}{}
    \pause = 2 + 2\times 9 - 9^2
    \pause = -61
  \]
\end{resolution}

%------------------------------------------------------------------------------%
\begin{resolution}{Segundo segmento}
  Segmento \(x=0\), \(0\leq y\leq 9\)

  \pause
  Função de uma variável, $y$, no intervalo fechado \(0\leq y\leq 9\)
  \[
    \pause h(y)
    \pause = f(0, y)
    \pause = \left(2 + 2x + 2y - x^2 - y^2\right)\evalat{x=0}{}
    \pause = 2 + 2y - y^2
  \]
  \pause
  Pontos críticos de $h$ no interior do intervalo \([0, 9]\)
  \vspace{-0.7\baselineskip}
  \begin{columns}[t]
  \column{0.4\linewidth}
    \pause
    \begin{align*}
      \onslide<+->{h'(y)}
      \onslide<+->{& = \frac{d}{dy}\left(2 + 2y - y^2\right) \\}
      \onslide<+->{& = 2 - 2y}
    \end{align*}
  \column{0.4\linewidth}
    \begin{align*}
      \onslide<+->{h'(y)  & = 0 \\}
      \onslide<+->{2 - 2y & = 0 \\}
      \onslide<+->{y      & = 1}
    \end{align*}
  \end{columns}
\end{resolution}

%------------------------------------------------------------------------------%
\begin{resolution}{Valores da função}
  Ponto crítico e extremos do intervalo \(y=0\), \(y=0\) e \(y=9\)

  \pause
  Pontos no plano
  \((0, 0)\), \((0, 1)\), \((0, 9)\)

  \pause
  Valores da função
  \[
    \pause
    f(0, 0) = 2
    \qquad\text{já calculado}
  \]
  \[
    \pause f(0, 1)
    \pause = h(1)
    \pause = \left(2 + 2y - y^2\right) \evalat{1}{}
    \pause = 3
  \]
  \[
    \pause f(0, 9)
    \pause = h(9)
    \pause = \left(2 + 2y - y^2\right) \evalat{9}{}
    \pause = 2 + 2\times 9 - 9^2
    \pause = -61
  \]
\end{resolution}

%------------------------------------------------------------------------------%
\begin{resolution}{Terceiro segmento}
  Segmento \(y=9-x\), \(0\leq x\leq 9\)

  \pause
  Função de uma variável, $x$, no intervalo fechado \(0\leq y\leq 9\)
  \begin{align*}
    \onslide<+->{p(x) & = f(x, 9-x)                        \\}
    \onslide<+->{& = 2 + 2x + 2(9-x) - x^2 - (9-x)^2       \\}
    \onslide<+->{& = 2 + 18 - 81 + 2x -2x +18x -x^2 - x^2  \\}
    \onslide<+->{& = -61 + 18x -2x^2}
  \end{align*}
\end{resolution}

%------------------------------------------------------------------------------%
\begin{resolution}{Terceiro segmento}
  \onslide<+->{Pontos críticos de $p$ no interior do intervalo $[0, 9]$}
  \vspace{-0.7\baselineskip}
  \begin{columns}[t]
  \column{0.5\linewidth}
    \begin{align*}
      \onslide<+->{p'(x)}
      \onslide<+->{& = \frac{d}{dx}\left(-61 + 18x -2x^2\right) \\}
      \onslide<+->{& = 18 - 4x}
    \end{align*}
  \column{0.5\linewidth}
    \begin{align*}
      \onslide<+->{p'(x)   & = 0           \\}
      \onslide<+->{18 - 4x & = 0           \\}
      \onslide<+->{9 - 2x  & = 0           \\}
      \onslide<+->{2x      & = 9           \\}
      \onslide<+->{x       & = \frac{9}{2}}
    \end{align*}
  \end{columns}
\end{resolution}

%------------------------------------------------------------------------------%
\begin{resolution}{Solução}
  Ponto crítico e extremos do intervalo \(x=0\), \(x=\frac{9}{2}\) e \(x=9\)

  \pause
  Para encontrar os pontos no plano usamos a relação \(y=9-x\)
  \vspace{0.5\baselineskip}

  \pause
  Quando \(x=0\),
  \pause
  \(
    \quad\, y = 9-0 = 9
  \)

  \pause
  Quando \(x=\frac{9}{2}\),
  \pause
  \(
    \quad y = 9-\frac{9}{2} = \frac{9}{2}
  \)

  \pause
  Quando \(x=9\),
  \pause
  \(
    \quad\, y = 9-9 = 0
  \)

  \pause
  \vspace{0.5\baselineskip}
  Pontos no plano
  \((0, 9)\), \(\left(\frac{9}{2}, \frac{9}{2}\right)\), \((9, 0)\)
\end{resolution}

%------------------------------------------------------------------------------%
\begin{resolution}{Valores da função}
  \begin{columns}
  \column{30mm}
    \onslide<+->{Já calculados}
    \begin{align*}
      \onslide<+->{f(0, 9) & = -61 \\}
      \onslide<+->{f(9, 0) & = -61}
    \end{align*}
  \column{80mm}
    \begin{align*}
      \onslide<+->{f\left(\frac{9}{2}, \frac{9}{2}\right)}
      \onslide<+->{& = p\left(\frac{9}{2}\right)                            \\}
      \onslide<+->{& = \left(-61 + 18x -2x^2\right)\evalat{\sfrac{9}{2}}{}  \\}
      \onslide<+->{& = -61 + 18\frac{9}{2} -2\left(\frac{9}{2}\right)^2     \\}
      \onslide<+->{& = \frac{-41}{2}}
    \end{align*}
  \end{columns}
\end{resolution}

%------------------------------------------------------------------------------%
\begin{resolution}{Solução}
  Máximo absoluto de $f$ em $R$ é 4 e ocorre no ponto \((1, 1)\)

  \vspace{\baselineskip}
  Mínimo absoluto de $f$ em $R$ é $-61$ e ocorre nos pontos \((9, 0)\) e \((0, 9)\)
\end{resolution}

%------------------------------------------------------------------------------%
